\section{Introduction}

Thrust vector control (TVC) is a fundamental technology in modern rocketry. It enables in-flight
trajectory correction and attitude stabilization by redirecting the thrust generated by a solid or
liquid powered rocket engine on a 2-axis gimbal. Since its development, TVC has played a crucial role in major space launch vehicles, from the Saturn and Space Shuttle programs through to modern vehicles such as SpaceX's Falcon 9 and NASA's Space Launch System (SLS) \cite{nasa_tvc_sls}. TVC allows the vectoring and angling of the thrust-producing component in order to counteract aerodynamic instabilities and initial launch trajectory errors.

While TVC is a core part of large-scale rockets, it remains niche and underdeveloped in model rocketry due to its high complexity and the manufacturing precision required, and most model rockets have relied on traditional, passive, aerodynamic stabilization using fins instead. The most notable exception is the hobbyist work of Joe Barnard's BPS.space project, which has demonstrated thrust-vectored model rockets capable of powered, controlled flight and even propulsive landing \cite{bpsspace}. This paper builds on that precedent by documenting, in full mathematical and software detail, a TVC system intended specifically as an educational and reproducible reference for other students.

The system described here is a compact, low-cost, and fully functional thrust vector control package, integrating a gimbal mount for thrust deflection, a high-speed PID control loop for attitude correction, and a Kalman filter for state estimation during powered flight. The vehicle operates entirely autonomously, from launch until
landing, transitioning through pre-determined flight states based on real-time sensor data.

The ultimate objective of this research is to demonstrate that TVC can be successfully adapted to
small-scale rocketry, and thereby increase the accessibility of controllable, reusable, and
autonomously governed model rockets to students and hobbyists. The paper also provides detailed
engineering insight into system design, flight software development, and simulation-based validation.
